\chapter{Lesson 4: Family and Honorifics}

Talking
about your family means talking about people older than you, and Korean grammar requires
you to show respect to them \emph{even when they are not present}. The tools are the
honorific infix \ko{\til (으)시\til}, special honorific verbs and nouns, and the
honorific particles \ko{께서} and \ko{께}.

\section{Family vocabulary}

Kinship terms for siblings depend on the \emph{speaker's} gender as well as the sibling's.

\begin{center}
\begin{tikzpicture}[every node/.style={font=\small, align=center},
    p/.style={draw, rounded corners, minimum width=2.1cm, minimum height=0.95cm, fill=blue!6},
    me/.style={draw, rounded corners, minimum width=2.1cm, minimum height=0.95cm, fill=orange!25, thick}]
  \node[p] (gf) at (-1.3,4.2) {\ko{할아버지}\\ \scriptsize harabeoji};
  \node[p] (gm) at (1.3,4.2) {\ko{할머니}\\ \scriptsize halmeoni};
  \draw (gf) -- (gm);
  \node[p] (uncle) at (-4.2,2.4) {\ko{삼촌 / 고모}\\ \scriptsize father's siblings};
  \node[p] (fa) at (-1.3,2.4) {\ko{아버지}\\ \scriptsize abeoji};
  \node[p] (mo) at (1.3,2.4) {\ko{어머니}\\ \scriptsize eomeoni};
  \node[p] (imo) at (4.2,2.4) {\ko{이모 / 외삼촌}\\ \scriptsize mother's siblings};
  \draw (fa) -- (mo);
  \draw (0,4.2) -- (0,3.3) -- (-4.2,3.3) -- (uncle.north);
  \draw (0,3.3) -- (-1.3,3.3) -- (fa.north);
  \draw[dashed] (mo.east) -- (imo.west);
  \node[p] (ob) at (-4.5,0.3) {\ko{형} / \ko{오빠}\\ \scriptsize older brother};
  \node[p] (os) at (-2.2,0.3) {\ko{누나} / \ko{언니}\\ \scriptsize older sister};
  \node[me] (me) at (0,0.3) {\ko{나} / \ko{저}\\ \scriptsize me};
  \node[p] (yb) at (2.2,0.3) {\ko{남동생}\\ \scriptsize younger brother};
  \node[p] (ys) at (4.5,0.3) {\ko{여동생}\\ \scriptsize younger sister};
  \draw (0,2.4) -- (0,1.3) -- (-4.5,1.3) -- (ob.north);
  \draw (-2.2,1.3) -- (os.north); \draw (0,1.3) -- (me.north);
  \draw (0,1.3) -- (4.5,1.3) -- (ys.north); \draw (2.2,1.3) -- (yb.north);
  \node[font=\scriptsize] at (0,-0.75) {$\longleftarrow$ older \hspace{4cm} younger $\longrightarrow$};
  \node[font=\scriptsize, align=left, anchor=west] at (-5.6,-1.5) {\ko{형}/\ko{오빠} and \ko{누나}/\ko{언니}: the first form is used by a male speaker, the second by a female speaker.};
\end{tikzpicture}
\end{center}

\begin{table}[H]
\centering
\begin{tabular}{llll}
\toprule
Term & RR & Meaning & Used by \\
\midrule
\ko{형} & hyeong & older brother & males \\
\ko{오빠} & oppa & older brother & females \\
\ko{누나} & nuna & older sister & males \\
\ko{언니} & eonni & older sister & females \\
\ko{동생} & dongsaeng & younger sibling & everyone \\
\ko{부모님} & bumonim & parents (honorific) & everyone \\
\ko{형제} & hyeongje & siblings & everyone \\
\ko{외할머니} & oehalmeoni & maternal grandmother & everyone \\
\ko{사촌} & sachon & cousin & everyone \\
\ko{남편 / 아내} & nampyeon / anae & husband / wife & everyone \\
\ko{아들 / 딸} & adeul / ttal & son / daughter & everyone \\
\bottomrule
\end{tabular}
\end{table}

\begin{remark}
The prefix \ko{외} oe- (``outside'') marks relatives on the mother's side:
\ko{외할아버지}, \ko{외할머니}, \ko{외삼촌}. The father's side is unmarked, a trace of the
traditional patrilineal family. \ko{아빠} appa and \ko{엄마} eomma are the affectionate
``dad'' and ``mom''; adults talking to non-family usually say \ko{아버지}, \ko{어머니}, or
\ko{저희 부모님} (my parents).
\end{remark}

\begin{example}[Describing your family]
Notice \ko{우리} / \ko{저희} (``our'') instead of ``my'' for family members:
\begin{itemize}
  \item \ko{우리 가족은 다섯 명이에요.} Uri gajogeun daseot myeongieyo. (There are five
        people in my family.)
  \item \ko{형이 한 명 있고 여동생이 한 명 있어요.} Hyeongi han myeong itgo yeodongsaengi
        han myeong isseoyo. (I have one older brother and one younger sister.)
  \item \ko{저는 형제가 없어요. 외동이에요.} Jeoneun hyeongjega eopseoyo. Oedongieyo. (I
        don't have any siblings. I'm an only child.)
\end{itemize}
\end{example}

\section{Grammar 1: The subject honorific \ko{\til (으)시\til}}

When the \emph{subject} of a verb is someone you should respect (a parent, grandparent,
teacher, customer, or a stranger who is older), insert \ko{\til (으)시\til} after the
stem.

\begin{itemize}
  \item Stem ends in a vowel: \ko{시}. \ko{가다} $\to$ \ko{가시다}.
  \item Stem ends in a consonant: \ko{으시}. \ko{읽다} $\to$ \ko{읽으시다}.
  \item Stem ends in \ko{ㄹ}: the \ko{ㄹ} drops and \ko{시} is added. \ko{살다} $\to$
        \ko{사시다}, \ko{알다} $\to$ \ko{아시다}.
\end{itemize}

In the polite style, \ko{시} + \ko{어요} contracts to \ko{세요}, and in the past
\ko{시} + \ko{었어요} gives \ko{셨어요}.

\begin{table}[H]
\centering
\begin{tabular}{lllll}
\toprule
Dictionary & Plain polite & Honorific present & Honorific past & Meaning \\
\midrule
\ko{가다} & \ko{가요} & \ko{가세요} & \ko{가셨어요} & go \\
\ko{오다} & \ko{와요} & \ko{오세요} & \ko{오셨어요} & come \\
\ko{읽다} & \ko{읽어요} & \ko{읽으세요} & \ko{읽으셨어요} & read \\
\ko{좋아하다} & \ko{좋아해요} & \ko{좋아하세요} & \ko{좋아하셨어요} & like \\
\ko{살다} & \ko{살아요} & \ko{사세요} & \ko{사셨어요} & live \\
\ko{바쁘다} & \ko{바빠요} & \ko{바쁘세요} & \ko{바쁘셨어요} & be busy \\
\ko{이다} & \ko{이에요/예요} & \ko{이세요/세요} & \ko{이셨어요/셨어요} & be \\
\bottomrule
\end{tabular}
\end{table}

\begin{important}[Who gets the honorific?]
The honorific shows respect to the \emph{subject}, never to yourself. Asking a teacher
\ko{어디 가세요?} (where are you going?) is correct; answering \ko{*도서관에 가세요} is
wrong --- you must say \ko{도서관에 가요}. Also, \ko{\til (으)세요} is the same form as the
polite command (``please do''), so \ko{여기 앉으세요} means ``please sit here''. Context
decides.
\end{important}

\section{Grammar 2: Special honorific verbs}

A handful of very common verbs have a completely different honorific form instead of
taking \ko{\til (으)시\til}.

\begin{table}[H]
\centering
\begin{tabular}{llllll}
\toprule
Plain & Meaning & Honorific & RR & Present & Past \\
\midrule
\ko{먹다}, \ko{마시다} & eat, drink & \ko{드시다} & deusida & \ko{드세요} & \ko{드셨어요} \\
\ko{자다} & sleep & \ko{주무시다} & jumusida & \ko{주무세요} & \ko{주무셨어요} \\
\ko{있다} & be (somewhere) & \ko{계시다} & gyesida & \ko{계세요} & \ko{계셨어요} \\
\ko{있다} & have & \ko{있으시다} & isseusida & \ko{있으세요} & \ko{있으셨어요} \\
\ko{말하다} & speak & \ko{말씀하시다} & malsseumhasida & \ko{말씀하세요} & \ko{말씀하셨어요} \\
\ko{아프다} & be sick & \ko{편찮으시다} & pyeonchaneusida & \ko{편찮으세요} & \ko{편찮으셨어요} \\
\ko{죽다} & die & \ko{돌아가시다} & doragasida & \ko{돌아가세요} & \ko{돌아가셨어요} \\
\bottomrule
\end{tabular}
\end{table}

\begin{remark}
\ko{있다} splits in two. When a respected person \emph{is somewhere}, use \ko{계시다}:
\ko{선생님은 사무실에 계세요} (the teacher is in the office). When they \emph{have}
something, use \ko{있으시다}: \ko{선생님은 시간이 있으세요?} (do you have time,
professor?). \ko{안녕히 계세요} (goodbye, to someone staying) literally means ``stay in
peace''. The more traditional \ko{잡수시다} japsusida also means ``eat'' and is used for
the very elderly.
\end{remark}

\section{Grammar 3: Honorific nouns and particles}

\begin{table}[H]
\centering
\begin{tabular}{lllll}
\toprule
Plain & Honorific & RR & Meaning & Typical phrase \\
\midrule
\ko{집} & \ko{댁} & daek & home & \ko{댁이 어디세요?} \\
\ko{나이} & \ko{연세} & yeonse & age & \ko{연세가 어떻게 되세요?} \\
\ko{이름} & \ko{성함} & seongham & name & \ko{성함이 어떻게 되세요?} \\
\ko{밥} & \ko{진지} & jinji & meal & \ko{진지 드셨어요?} \\
\ko{말} & \ko{말씀} & malsseum & words, speech & \ko{말씀 많이 들었어요.} \\
\ko{생일} & \ko{생신} & saengsin & birthday & \ko{생신 축하드려요.} \\
\ko{사람}, \ko{명} & \ko{분} & bun & person & \ko{몇 분이세요?} \\
\bottomrule
\end{tabular}
\end{table}

The particles change too:
\begin{itemize}
  \item subject \ko{이/가} $\to$ \ko{께서} kkeseo; topic \ko{은/는} $\to$ \ko{께서는}.
  \item ``to (a person)'' \ko{에게/한테} $\to$ \ko{께} kke.
\end{itemize}
And on the speaker's own side there are \textbf{humble} forms: \ko{저} (I),
\ko{저희} (we), \ko{드리다} deurida (give to someone higher, replacing \ko{주다}),
\ko{여쭤보다} yeojjwoboda (ask someone higher, replacing \ko{물어보다}).

\begin{center}
\begin{tikzpicture}[every node/.style={font=\small},
    w/.style={draw, rounded corners, minimum height=0.8cm, fill=gray!8},
    h/.style={draw, rounded corners, minimum height=0.8cm, fill=red!12}]
  \node[w] (a) at (0,1.4) {\ko{동생이}};
  \node[w] (b) at (2.4,1.4) {\ko{밥을}};
  \node[w] (c) at (4.6,1.4) {\ko{먹어요.}};
  \node[h] (d) at (0,0) {\ko{할머니께서}};
  \node[h] (e) at (2.4,0) {\ko{진지를}};
  \node[h] (f) at (4.6,0) {\ko{드세요.}};
  \draw[-{Stealth}, gray] (a) -- (d); \draw[-{Stealth}, gray] (b) -- (e); \draw[-{Stealth}, gray] (c) -- (f);
  \node[anchor=west, font=\scriptsize] at (6.0,1.4) {My younger sibling eats.};
  \node[anchor=west, font=\scriptsize] at (6.0,0) {My grandmother eats.};
  \node[font=\scriptsize] at (0,-0.65) {particle};
  \node[font=\scriptsize] at (2.4,-0.65) {noun};
  \node[font=\scriptsize] at (4.6,-0.65) {verb};
\end{tikzpicture}
\end{center}

\begin{example}[Honorific sentences]
All three layers (particle, noun, verb) shift together:
\begin{itemize}
  \item \ko{할아버지께서 신문을 읽으세요.} Harabeojikkeseo sinmuneul ilgeuseyo. (Grandpa
        is reading the newspaper.)
  \item \ko{할머니께서는 지금 주무세요.} Halmeonikkeseoneun jigeum jumuseyo. (Grandma is
        sleeping now.)
  \item \ko{아버지께서는 회사에 계세요.} Abeojikkeseoneun hoesae gyeseyo. (My father is
        at the office.)
  \item \ko{어머니께 꽃을 드렸어요.} Eomeonikke kkocheul deuryeosseoyo. (I gave flowers to
        my mother.)
  \item \ko{할아버지는 작년에 돌아가셨어요.} Harabeojineun jangnyeone doragasyeosseoyo.
        (My grandfather passed away last year.)
  \item \ko{선생님, 주말에 뭐 하셨어요?} Seonsaengnim, jumare mwo hasyeosseoyo?
        (Professor, what did you do on the weekend?)
\end{itemize}
\end{example}

\begin{important}[Two independent axes]
Korean politeness has two separate dimensions, which students mix up:
\begin{itemize}
  \item \textbf{Speech level} (the sentence ending, \ko{\til 요} or not) shows respect to
        the \emph{listener}.
  \item \textbf{Honorifics} (\ko{\til 시\til}, \ko{께서}, \ko{드시다}) show respect to the
        \emph{subject} being talked about.
\end{itemize}
So talking to a friend about your grandmother gives \ko{할머니가 주무셔} (casual ending,
honorific verb), while talking to a teacher about your little brother gives
\ko{동생이 자요} (polite ending, plain verb).
\end{important}

\section{Asking politely about someone's family}

\begin{example}[Polite questions]
Typical first-meeting questions to an older person:
\begin{itemize}
  \item \ko{가족이 몇 분이세요?} Gajogi myeot buniseyo? (How many are in your family?)
  \item \ko{부모님은 어디 사세요?} Bumonimeun eodi saseyo? (Where do your parents live?)
  \item \ko{아버님은 무슨 일을 하세요?} Abeonimeun museun ireul haseyo? (What does your
        father do?) --- \ko{아버님} and \ko{어머님} are used for \emph{someone else's}
        parents.
\end{itemize}
\end{example}

\begin{aside}
\textbf{Titles, age and family words outside the family.} Koreans rarely call anyone
older than them by their bare name. Family terms are extended to friends and
acquaintances: a male student calls a slightly older male friend \ko{형} and an older
female friend \ko{누나}; a female student uses \ko{오빠} and \ko{언니}. Restaurant staff
may be called \ko{이모} (aunt). At university, upperclassmen are \ko{선배} seonbae and
underclassmen are \ko{후배} hubae; even a one-year difference matters, and \ko{선배} are
traditionally expected to pay when they take \ko{후배} out to eat. With colleagues and
classmates of uncertain status, use full name + \ko{씨} ssi (\ko{김민수 씨}), never
surname + \ko{씨}, and avoid the dictionary word for ``you'' (\ko{당신}), which sounds
confrontational or overly intimate. Teachers are \ko{선생님}, and the suffix \ko{\til 님}
attaches to job titles: \ko{교수님} (professor), \ko{사장님} (company president).
\end{aside}

\section*{Review summary}
\begin{itemize}
  \item Sibling terms depend on your gender: \ko{형/누나} (male speaker),
        \ko{오빠/언니} (female speaker); \ko{동생} for younger siblings.
  \item Honorific \ko{\til (으)시\til}: \ko{가세요}, \ko{읽으세요}, \ko{가셨어요};
        \ko{ㄹ}-stems drop \ko{ㄹ} (\ko{사세요}).
  \item Special verbs: \ko{드시다, 주무시다, 계시다, 말씀하시다, 편찮으시다,
        돌아가시다}.
  \item Special nouns: \ko{댁, 연세, 성함, 진지, 말씀, 생신, 분}; particles \ko{께서},
        \ko{께}; humble \ko{저, 저희, 드리다}.
  \item Never honor yourself; speech level (listener) and honorifics (subject) are
        separate choices.
\end{itemize}
